\documentclass[12pt,english]{article}
\bibliographystyle{jf}
\usepackage[longnamesfirst]{natbib}
%\usepackage[T1]{fontenc}
\usepackage{fontenc}
\usepackage[utf8]{inputenc}
\usepackage[bottom]{footmisc}
\usepackage{tabularx}
\usepackage{booktabs}
\usepackage{amsmath,enumerate}
\usepackage{amsthm}
\usepackage{dsfont}
\usepackage{amssymb}
\usepackage{mathtools}
\usepackage{mathrsfs}
\usepackage{graphicx}
\usepackage{natbib}
\usepackage[labelfont=bf,justification=justified,font=small,skip=2pt,labelsep=period]{caption}
%\usepackage[labelfont=bf,justification=justified,labelsep=period]{caption}
\usepackage{subcaption}
\usepackage{setspace}
\usepackage{arydshln}
\usepackage{chngpage}
\usepackage{float,appendix}
\usepackage{afterpage}
\usepackage{verbatim}
\usepackage{indentfirst}
\usepackage[margin=0.7in]{geometry}
\usepackage[usenames, dvipsnames]{color}
\usepackage{hyperref}
\usepackage{adjustbox}
% \usepackage{longtable}
\usepackage{booktabs, makecell, longtable,ltxtable}
\usepackage{lipsum}
\usepackage{pdflscape}
\usepackage{siunitx}
\usepackage{chngcntr}
\usepackage{lipsum}
\usepackage{tikz}
\usepackage{pdflscape}
\usepackage[normalem]{ulem}
\usepackage{chngcntr,apptools}
\AtAppendix{\counterwithin{lemma}{section}}
\usepackage{siunitx}
\usepackage{tabularx}
\sisetup{table-format=-1.3, table-space-text-post={**}}
\usepackage{xr}
\usepackage{etoolbox}
\usetikzlibrary{positioning}
\newcommand\fnsep{\textsuperscript{,}}
\usepackage{tikz}
\usepackage{pdflscape}
\usepackage[normalem]{ulem}
\usepackage{chngcntr,apptools}
\AtAppendix{\counterwithin{lemma}{section}}
\usepackage{siunitx}
\usepackage{tabularx}
\usepackage{xr}
\usepackage{etoolbox}
\usetikzlibrary{positioning}

%% UCLA colors: 
% https://brand.ucla.edu/wp-content/uploads/ucla-color-palette-v10.pdf
%\definecolor{UCLA}{RGB}{50, 132, 191}   % UCLA Blue primary
%\definecolor{UCLA}{RGB}{52, 123, 173}   % UCLA Blue text only
\definecolor{UCLAblue}{RGB}{30, 75, 135} % secondary color  
\definecolor{UCLAgold}{RGB}{255, 232, 0} % UCLA Gold

\definecolor{usccardinal}{rgb}{0.6, 0.0, 0.0}
\definecolor{Matlabblue}{rgb}{0,0.447,0.741}
\definecolor{bleudefrance}{rgb}{0.19,0.55, 0.91}
\definecolor{cobalt}{rgb}{0.0, 0.28,0.67}
\definecolor{darkblue}{rgb}{0.0, 0.0,0.55}
\definecolor{darkcerulean}{rgb}{0.03,0.27, 0.49}
\definecolor{darkpowderblue}{rgb}{0.0,0.2, 0.6}
\definecolor{bleudefrance}{rgb}{0.19,0.55, 0.91}
%\usepackage[colorlinks=true,urlcolor=blue,linkcolor=blue,citecolor=blue]{hyperref}

\hypersetup{
  colorlinks,
  linkcolor={UCLAblue},
  citecolor={UCLAblue},
  urlcolor={UCLAblue},
}

\bibpunct{(}{)}{;}{a}{,}{,}
\newcommand{\E}{\mathbb{E}}
\newcommand{\var}{\mathrm{Var}}
\newcommand{\cov}{\mathrm{Cov}}
\setcounter{MaxMatrixCols}{20}
\newtheorem{theorem}{Theorem}
\newtheorem{lemma}{Lemma}
\newtheorem{defn}{Definition}
\newtheorem{prop}{Proposition}
\newtheorem{cor}{Corollary}
\urlstyle{same}
\usepackage{pgfplots} 
\pgfplotsset{compat=newest} 
\pgfplotsset{plot coordinates/math parser=false} 
% puts figures and tables at the end of the document
%\usepackage[nomarkers,nolists]{endfloat}               
% puts ONLY figures at the end of the document
%\usepackage[nomarkers,nolists,figuresonly]{endfloat}  
% allows more than one figure per page for endfloat
%\renewcommand{\efloatseparator}{\mbox{}}              
\usepackage{rotating}
% allows sideways tables and figures to be placed at the end of the document
%\DeclareDelayedFloatFlavor{sidewaystable}{table}
%\DeclareDelayedFloatFlavor{sidewaysfigure}{figure}

% \usepackage[nomarkers,nolists]{endfloat} 
% \renewcommand{\efloatseparator}{\mbox{}}              % allows more than one figure per page for endfloat

%\usepackage{mathpazo}
%\usepackage{mathpple}
%%%%%%%%%%%%%%%%%%%%%%%%%%%%%%%%
%\usepackage{palatino}
%%%%%%%%%%%%%%%%%%%%%%%%%%%%%%%%
\pgfplotsset{compat=newest} 
\pgfplotsset{plot coordinates/math parser=false} 
\newlength\figureheight 
\newlength\figurewidth 
\usepackage{mathtools}
\usetikzlibrary{calc}


%%%%%%%%%%%%%%%%%%%%%%%%%%
% Notes customization %
%%%%%%%%%%%%%%%%%%%%%%%%%%
\definecolor{webblue}{rgb}{0,0,0.6}
\definecolor{webgreen}{rgb}{0,.6,0}
\definecolor{webbrown}{rgb}{.6,0,0}


%- Commenting
\let\oldmarginpar\marginpar
\renewcommand\marginpar[1]{\-\oldmarginpar[\raggedleft\tiny #1]%
{\raggedright\tiny #1}}
\newcommand{\mknote}[1]{\marginpar{\textcolor{webblue}{\tiny{Mahyar: #1}}}}
\newcommand{\dsnote}[1]{\marginpar{\textcolor{webgreen}{\tiny{Dejanir: #1}}}}
